\documentclass{article}
\usepackage{amssymb,amsmath,blindtext,listings,graphicx,siunitx, mathrsfs, url}
\title{Special Relativity Equation Cheat Sheat}
\date{2026 May}
\author{Sean Brennan \\ www.zettix.com \\ github.com/zettix/SpecialRelativityEquationsCheatSheet}
\begin{document}
\maketitle

\section{Introduction}
This is a collection of equations to reference when learning SR on Youtube.
Stanford has a class available here: Special Relativity, Stanford.
Taught by Leonard Susskind.
\url{https://www.youtube.com/watch?v=Z8ace3NbXAA&list=PLD9DDFBDC338226CA}

\section{The Equations}

Speed of light

\begin{equation}
\begin{split}
c &\approx 300,000,000\:\si{\meter/\second} = 3 \times 10^8\:\si{\meter/\second}\\
&= 299,792,458\:\si{\meter/\second} 
\end{split}
\end{equation}

Proper Time (Invariant)

\begin{equation}
\tau^2 = t^2 - x^2
\end{equation}

\begin{equation}
\Delta \tau^2 = \Delta t^2 - \Delta x^2
\end{equation}

Lorenz Transformations

\begin{equation}
t' = \frac{t - vx}{\sqrt{1 - v^2}}
\end{equation}

\begin{equation}
t = \frac{t' - vx'}{\sqrt{1 - v^2}}
\end{equation}

\begin{equation}
x' = \frac{x - vt}{\sqrt{1 - v^2}}
\end{equation}

\begin{equation}
x = \frac{x' - vt'}{\sqrt{1 - v^2}}
\end{equation}

Newtonian Momentum

\begin{equation}
p = mv
\end{equation}

Newtonian Force

\begin{equation}
f = ma = mv^2 = pv
\end{equation}

4-Space: with time, $X$ has $\mu$ superscript.
\begin{equation}
X^\mu = (t, x, y, z)
\end{equation}

4-Space: without time, $X$ has $i$ superscript.
\begin{equation}
X^i = (x, y, z)
\end{equation}

4-Dimensional Velocity
\begin{equation}
\frac{\Delta X^\mu}{\Delta \tau}
\end{equation}

4-Dimensional Time Velocity
\begin{equation}
U_0 = U^0 = \frac{1}{\sqrt{1 - v^2}}
\end{equation}

4-Dimensional Space Velocity
\begin{equation}
\vec U = U^i = U_i = \frac{v^i}{\sqrt{1 - v^2}}
\end{equation}

Note
\begin{equation}
U_0^2 - \vec U^2 = 1
\end{equation}

Action Potental Sums
\begin{equation}
A = -m\sum{\Delta \tau}
\end{equation}

Action Potental Integral
\begin{equation}
\begin{split}
A  &= -m\int_a^b{\Delta \tau}\\
&= -m\int_a^b{\sqrt{dt^2 - dx^2}}\\
&= -m\int_a^b{\sqrt{1 - \frac{dx^2}{dt^2}} dt}\\
&= -m\int_a^b{\sqrt{1 - \dot x^{i^2}} dt}
\end{split}
\end{equation}

Lagrangian
\begin{equation}
\begin{split}
\mathscr{L} &= -m\sqrt{1 - \left( \dot x^{2} + \dot y^{2} + \dot z^{2} \right) }\\
&=-m\sqrt{1-v^2}
\end{split}
\end{equation}

Lagrangian's Lagrangian
\begin{equation}
\mathscr{L}_{\mathscr{L}} = \frac{m \dot x^2}{2}
\end{equation}
where
\begin{equation}
\dot x^2 = \sum{ \dot x_i^{2}}
\end{equation}

Continue the Action potential integral and get Newtonian Energy
\begin{equation}
\begin{split}
A  &= -m\int_a^b{\mathscr{L}}\\
&= -m\left( 1 - \frac{v^2}{2} \right)\\
&= mv^2
\end{split}
\end{equation}
igoring the constant $-m$ from the final value.

\section{Conservation of Energy and Momentum}

Lagrangian momentum

\begin{equation}
\begin{split}
P_x &= \frac{\partial\mathscr{L}}{\partial  \dot x}\\
&= \frac{m \dot x}{2 \sqrt{1 - v^2}}\\
&= \frac{m v_x}{\sqrt{1-v^{2}}}
\end{split}
\end{equation}

where
\begin{equation}
v^2 = x^2 + y^2 + z^2
\end{equation}

This looks like $U_x$, hence, relativistic momemtum uses the mass at rest and the relativistic space axes $x$, $y$, $z$. e.g. $x$:
\begin{equation}
P_x = \frac{m v_x}{\sqrt{1-v^2}}
\end{equation}

4-vector in terms of deltas $\Delta t$, $\Delta x$.

Momentum by component
\begin{equation}
\begin{split}
mU^x &= P_x\\
mU^y &= P_y\\
mU^z &= P_z\\
mU^0 &= P_0
\end{split}
\end{equation}

Conservation of momenta
\begin{equation}
\vec P_{initial} = \vec P_{final}
\end{equation}

or

\begin{equation}
\vec P_{initial} - \vec P_{final} = 0
\end{equation}

What is $P_0$?  Is it the particle energy?  Look at the Hamiltonian.
\begin{equation}
H = \sum_i{\dot X_i P^i - \mathscr{L}}
\end{equation}

By component, $P^i$
\begin{equation}
P^i = \frac{m \dot X^{i}}{\sqrt{1 - v^2}}
\end{equation}

By component, $\mathscr{L}$
\begin{equation}
\mathscr{L} = -m \sqrt{1 - v^{2}}
\end{equation}

Rewriting the Hamiltonian
\begin{equation}
H = \sum_i \frac{m \dot x_{i}^{2}}{\sqrt{1 - v^2}} + m \sqrt{1-v^2}
\end{equation}

using

\begin{equation}
v^2 = \sum_i \dot x_{i}^{2}
\end{equation}

we get and simplify
\begin{equation}
\begin{split}
H &=\frac{m v^2}{\sqrt{1 - v^2}} + m \sqrt{1-v^2} \\
&=\frac{m v^2}{\sqrt{1 - v^2}} + \frac{m \left( \sqrt{1-v^2} \right)^2}{\sqrt{1-v^2}}\\
&=\frac{m v^2 + m - m v^2}{\sqrt{1-v^2}}\\
&=\frac{m}{\sqrt{1-v^2}} = m U_0
\end{split}
\end{equation}

To calculate energy of a particle at rest, we use the Taylor expansion of
\begin{equation}
\frac{1}{\sqrt{1 - v^2}} = \left( 1 - v^2 \right)^{-\frac{1}{2}} \approx 1 + \frac{v^2}{2}
\end{equation}

And get energy

\begin{equation}
H = m + \frac{m v^2}{2}
\end{equation}

Note that at rest, velocity is zero so energy is simply $m$, however, the units of $H$, to be consistant with
$\frac{m v^2}{2}$ must be $\frac{g m^2}{s^2}$ hence $\frac{m v^2}{2}$ drops and fixing the units, rest energy is:
\begin{equation}
H_{rest} = mc^2
\end{equation}

\section{Massless Particles}

Clearly, $m=0$, $v=c$ causes energy, $\frac{m}{1 - v^2}$, to go to $\frac{0}{0}$ so instead we use momentum. Recall


\begin{equation}
U^{0^2} - U^{x^2} - U^{y^2} - U^{z^2} = 1
\end{equation}

Multiply by m

\begin{equation}
m^2U^{0^2} - m^2U^{x^2} - m^2U^{y^2} - m^2U^{z^2} = m^2
\end{equation}

Substitute energy and momentum
\begin{equation}
E^2 - P_x^2 - P_y^2 - P_z^2 = m^2
\end{equation}

To Energy:

\begin{equation}
E = \sqrt{P^2 + m^2}
\end{equation}

With units

\begin{equation}
E = \sqrt{P^2 c^2 + m^2 c^4}
\end{equation}

As massless:
\begin{equation}
E = c |P|
\end{equation}

\section{positronium}

An electon and an positron in orbit, and will decay into two photons, leaving in opposite directions.
Calculate the energy and momentum of the system before and after decay.  Initial mass in joules or grams.
Momentum is conserved.  At rest, zero.  Photons must go in opposite directions for momentum to be zero.  And equal momentum, to ensure zero momentum.

Start with
\begin{equation}
E_{rest} = m_{positronium} c^2
\end{equation}

This must equal the end state, twice the momenta of each photon:

\begin{equation}
m c^2 = 2 c |P|
\end{equation}

Rearrange
\begin{equation}
\frac{m c}{2} = |P|
\end{equation}

\section{Energy with Units}

\begin{equation}
E = \frac{m c^2}{\sqrt{1 - \frac{v^2}{c^2}}}
\end{equation}

\section{Lecture 4}

Fields. Discussion of field function from $t$ and $x$ to $\phi$.
\begin{equation}
\phi(t, x)
\end{equation}

The action

\begin{equation}
Action = \int_a^b \frac{m}{2}\left( \frac{d\phi}{dt} \right)^2 - V(\phi) dt
\end{equation}

Lagrangian

\begin{equation}
\mathscr{L} = \frac{m}{2} \dot \phi^{2} - V(\phi)
\end{equation}

Equations of motion given by

\begin{equation}
\frac{d}{dt} \: \frac{\partial \mathscr{L}}{\partial \dot \phi} =\\
\frac{\partial \mathscr{L}}{\partial \phi}
\end{equation}

Carry out computation to get Newton's equations:
\begin{equation}
m \ddot \phi = - \frac{\partial V}{\partial \phi}
\end{equation}

Discussion of least action from point $a$ to point $b$.  Cut time axis into
slices, the integral is a sum of the value of $phi$ at each slice.  To minimize,
take the derivative.


We are $Three\: plus\: one$ for three space dimensions and one time dimesion.

Hold fixed the boundary, wiggle up and down, until we minimize the action.

In practice, the Action Potential is

\begin{equation}
A  = \int \int \int \int dx\: dy\: dz\: dt
\end{equation}

But in shorthand
\begin{equation}
A  = \int d^4 x \: \mathscr{L} \left(\phi, \frac{\partial \phi}{\partial t},\\
\frac{\partial \phi}{\partial x}, \frac{\partial \phi}{\partial y},\\
\frac{\partial \phi}{\partial z}\right)
\end{equation}

Condense once
\begin{equation}
A  = \int d^4 x \: \mathscr{L} \left(\phi, \frac{\partial \phi}{\partial x^\mu} \right)
\end{equation}

And use
\begin{equation}
\begin{split}
\phi_t &= \frac{\partial \phi}{\partial t}\\
\phi_x &= \frac{\partial \phi}{\partial x}\\
\phi_y &= \frac{\partial \phi}{\partial y}\\
\phi_z &= \frac{\partial \phi}{\partial z}
\end{split}
\end{equation}

To condense notation further, rather than use $\dot \phi$

How to handle a field, which is a surface?   Some other notes describe the
Euler Lagrange equations, but here we have:

\begin{equation}
\frac{d}{dt} \: \frac{\partial \mathscr{L}}{\partial \dot \phi} -\\
\frac{\partial \mathscr{L}}{\partial \phi} = 0
\end{equation}

Or starting here:
\begin{equation}
\frac{d}{dt} \: \frac{\partial \mathscr{L}}{\partial \frac{d \phi}{dt}} -\\
\frac{\partial \mathscr{L}}{\partial \phi} = 0
\end{equation}



\end{document}

